\section{Conclusion}\label{sec:conclusion}
This study audited covariate adjustment in patient-level classification of the TUEP label from first-session routine EEG. When the EEG carried no label information, adjustment created up to 11.9 pp of discrimination in nonlinear classifiers on real EEG features (plasmode), and a reference fitted on training controls leaked even with a linear classifier.
The absolute decision rule for the feature-shuffle test rarely raised false alarms and detected leaks above an AUC of 55, whereas the paired rule, with 20--30 replicates, was often undecided and, for raw tree pipelines, declared leakage in 32\%--48\% of null draws. Rank-normal mitigation reduced the leak in tree classifiers but not enough in a sequence model. Within patients, the leak was not shown to grow with the correlation of the windows, and, in one simulated cell, aggregating to patient means before adjusting helped only at low correlation.
The audited fusion pipeline was not shown to exceed the acquisition-metadata baseline or a model without adjustment; discrimination persisted within matched strata, but it was not shown to be specific to F2 and cannot yet be separated from antiseizure medication and report-derived labels.
The evidence comes from simulations and one retrospective archive ($n = 200$) after 49 rounds of analysis and does not support clinical claims.

\section*{CRediT authorship contribution statement}
\ISI{Name: CRediT role}
